\chapter{Geometry of the Cone Generators: Planes, Angles, and Growth}
\label{sec:kegelstrahlen}

\section{Overview}

The integer points of the sequence $k(n)$ lie on discrete
generators of the cone. Every two adjacent generators span a
triangular face whose position in space is completely determined
by the cone parameters. In this chapter we derive the plane equation,
the tilt angle, the angles between the generators, and the
growth rate of the radii. As a concrete example, the triangular
area at $n = 19$ is computed.

% ─────────────────────────────────────────────────────────────────────────────
\section{Fundamentals: Generators of the Cone}
% ─────────────────────────────────────────────────────────────────────────────

The cone is defined by the equation
\[
  z = -\frac{3}{4} + \frac{\pi}{2}\,r
\]
with apex at the point $P_0 = (0, 0, -\tfrac{3}{4})$.
A generator at azimuthal angle $\varphi$ runs from $P_0$ in the direction
\[
  \vec{v}(\varphi)
  = \begin{pmatrix}\cos\varphi\\ \sin\varphi\\ \pi/2\end{pmatrix}.
\]
The factor $\pi/2$ in the $z$-component follows directly from the
cone slope $dz/dr = \pi/2$.

The length of this direction vector (per unit $\Delta r = 1$) is:
\[
  |\vec{v}| = \sqrt{1^2 + \left(\frac{\pi}{2}\right)^2}
            = \sqrt{1 + \frac{\pi^2}{4}}
            \approx 1.86209589.
\]

\subsection*{Integer Points and Their Radii}

At height $z = n$ the point lies on the generator at radius
\[
  r_n = \frac{z + \tfrac{3}{4}}{\pi/2}
      = \frac{2}{\pi}\!\left(n + \frac{3}{4}\right)
      = \frac{2n}{\pi} + \frac{3}{2\pi}.
\]
The radius grows \emph{linearly} in $n$ at the rate

\begin{equation}\label{eq:wachstumsrate}
  \frac{dr_n}{dn} = \frac{2}{\pi} \approx 0.63661977.
\end{equation}

The following table shows the radii for the first integer
indices $n \equiv 1, 5 \pmod{6}$:

\begin{table}[H]
\centering
\renewcommand{\arraystretch}{1.2}
\caption{Radii $r_n = \tfrac{2}{\pi}n + \tfrac{3}{2\pi}$ of the integer
  sequence values on the cone.}
\begin{tabular}{ccc}
\toprule
$n$ & $r_n = \tfrac{2n}{\pi}+\tfrac{3}{2\pi}$ & $r_n$ (numerical) \\
\midrule
 1  & $\tfrac{2}{\pi}+\tfrac{3}{2\pi} = \tfrac{7}{2\pi}$  &  1.11408460 \\
 5  & $\tfrac{10}{\pi}+\tfrac{3}{2\pi} = \tfrac{23}{2\pi}$ &  3.66056369 \\
 7  & $\tfrac{14}{\pi}+\tfrac{3}{2\pi} = \tfrac{31}{2\pi}$ &  4.93380324 \\
11  & $\tfrac{22}{\pi}+\tfrac{3}{2\pi} = \tfrac{47}{2\pi}$ &  7.48028233 \\
13  & $\tfrac{26}{\pi}+\tfrac{3}{2\pi} = \tfrac{55}{2\pi}$ &  8.75352187 \\
17  & $\tfrac{34}{\pi}+\tfrac{3}{2\pi} = \tfrac{71}{2\pi}$ & 11.30000096 \\
19  & $\tfrac{38}{\pi}+\tfrac{3}{2\pi} = \tfrac{79}{2\pi}$ & 12.57324050 \\
\bottomrule
\end{tabular}
\end{table}

\begin{remark}[Arithmetic of the radius column]
\label{rem:radius_arithmetic}
The table encodes three layers of arithmetic beyond the linear growth.

(i) \emph{Integer circumferences.} Combining the two fractions gives
$r_n = (4n+3)/(2\pi) = Q/(2\pi)$, hence
\begin{equation}\label{eq:circumference_Q}
  2\pi r_n = 4n+3 = Q = \sqrt{48\,k(n)+1}:
\end{equation}
the circumference of the cone at each integer point is an integer,
namely the universal quantity $Q$ (Section~\ref{subsec:Q}). The
numerators $7, 23, 31, 47, \ldots$ in the middle column are these
circumferences. Modulo $12$ they separate the two residue classes:
$Q \equiv 7$ for $n \equiv 1 \pmod 6$ and $Q \equiv 11$ for
$n \equiv 5 \pmod 6$ --- the circumference reads off the class
membership.

(ii) \emph{OEIS anchors.} The index sequence $n = 1, 5, 7, 11, \ldots$
is A007310 in~\cite{OEIS}; the numerators $2n = 2, 10, 14, 22, \ldots$
of the linear term $2n/\pi$ run through A091999, the numbers congruent
to $2$ or $10$ modulo $12$ --- equivalently, the even numbers divisible
by neither $3$ nor $4$. Doubling maps the classes $\{1,5\} \bmod 6$
bijectively onto $\{2,10\} \bmod 12$; the twin pattern of gaps $(4,2)$
becomes $(8,4)$. Each element $2n$ with $n > 1$ is the sum of exactly
four consecutive integers, $2n = m + (m{+}1) + (m{+}2) + (m{+}3)$ with
$m = (n-3)/2$, and the starting values again track the classes:
$m \equiv 1 \pmod 3$ for $n \equiv 5 \pmod 6$ and $m \equiv 2 \pmod 3$
for $n \equiv 1 \pmod 6$.

(iii) \emph{An infinite product evaluates the angular quantum.} With
$a(j) = 6j - 3 + (-1)^j$ enumerating A091999, one has the
identity~\cite{OEIS}
\begin{equation}\label{eq:product_sin}
  \prod_{j=1}^{\infty} \Bigl(1 + \frac{(-1)^j}{a(j)}\Bigr)
  = 2\sin\frac{\pi}{12} = \frac{\sqrt{6}-\sqrt{2}}{2}
  = 0.51763809\ldots,
\end{equation}
confirmed here numerically to seven digits with $2 \cdot 10^6$
factors. The angle $\pi/12$ is precisely the angle per $Q$-unit of the
parametrisation, $d\theta = (\pi/12)\,dQ$
(Theorem~\ref{thm:q-arc-length}), and by the reflection formula
$\Gamma(\tfrac{1}{12})\,\Gamma(\tfrac{11}{12}) = \pi/\sin(\pi/12)$ the
identity is the twelfth-argument sibling of
Remark~\ref{rem:gamma_reflection}. The product evaluation itself is
classical; the connection to the constants of the present construction
is an observation of this work.
\end{remark}

\begin{bemerkung}[Linear Growth Rate]
The radial difference between two consecutive generators
is constant:
\[
  \Delta r_{\text{Minor}} = \frac{2}{\pi}\cdot 2 = \frac{4}{\pi}
  \approx 1.27323954,
  \qquad
  \Delta r_{\text{Major}} = \frac{2}{\pi}\cdot 4 = \frac{8}{\pi}
  \approx 2.54647909.
\]
This is a direct consequence of the linearity~\eqref{eq:wachstumsrate}.
\end{bemerkung}

% ─────────────────────────────────────────────────────────────────────────────
\section{Angular Separations of the Generators in the \texorpdfstring{$xy$}{xy}-Plane}
% ─────────────────────────────────────────────────────────────────────────────

The azimuthal angles of the integer points follow from
$\theta_n = \tfrac{\pi}{3}n + \tfrac{\pi}{4}$. Consecutive
points differ by:

\begin{align*}
  \Delta\varphi_{\text{Minor}} &= \frac{\pi}{3}\cdot 2 = \frac{2\pi}{3}
    = 120^\circ, \\
  \Delta\varphi_{\text{Major}} &= \frac{\pi}{3}\cdot 4 = \frac{4\pi}{3}
    = 240^\circ.
\end{align*}

Together they yield a complete revolution:
$\Delta\varphi_{\text{Minor}} + \Delta\varphi_{\text{Major}} = 2\pi$.

% ─────────────────────────────────────────────────────────────────────────────
\section{Angle Between Two Generators in Space}
% ─────────────────────────────────────────────────────────────────────────────

Two generators with azimuthal difference $\Delta\varphi$ have the
direction vectors
\[
  \vec{v}_1 = \begin{pmatrix}1\\0\\\pi/2\end{pmatrix},
  \qquad
  \vec{v}_2 = \begin{pmatrix}\cos\Delta\varphi\\\sin\Delta\varphi\\\pi/2\end{pmatrix}.
\]

The dot product reads:
\[
  \vec{v}_1 \cdot \vec{v}_2
  = \cos\Delta\varphi + \left(\frac{\pi}{2}\right)^2
  = \cos\Delta\varphi + \frac{\pi^2}{4}.
\]

Since $|\vec{v}_1| = |\vec{v}_2| = \sqrt{1+\pi^2/4}$, we obtain:
\[
  \cos\theta
  = \frac{\cos\Delta\varphi + \pi^2/4}{1 + \pi^2/4}.
\]

For minor and major steps, $\cos(120^\circ) = \cos(240^\circ) = -\tfrac{1}{2}$,
hence:
\[
  \cos\theta
  = \frac{-\tfrac{1}{2} + \tfrac{\pi^2}{4}}{1 + \tfrac{\pi^2}{4}}
  = \frac{\pi^2 - 2}{\pi^2 + 4}
  \approx 0.56739934.
\]

\begin{satz}[Angle Between the Generators]
The spatial angle between two consecutive generators
(both minor and major steps) is:
\begin{equation}
  \theta = \arccos\!\left(\frac{\pi^2 - 2}{\pi^2 + 4}\right)
  \approx 55.43092771^\circ.
\end{equation}
\end{satz}

\begin{bemerkung}
Although the azimuthal angles $120^\circ$ and $240^\circ$ are different, the
spatial angle between the generators is the same in both cases.
The reason: $\cos(120^\circ) = \cos(240^\circ) = -\tfrac{1}{2}$, so that
the $z$-component $\pi/2$ has the same effect in both cases,
and the different opening angles in the plane lead to the same
spatial angle.
\end{bemerkung}

% ─────────────────────────────────────────────────────────────────────────────
\section{Plane Equation and Tilt Angle of the Triangular Face}
% ─────────────────────────────────────────────────────────────────────────────

Two generators span a triangular face. Its
normal vector is obtained via the cross product:
\[
  \vec{n}
  = \vec{v}_1 \times \vec{v}_2
  = \begin{vmatrix}\vec{e}_x & \vec{e}_y & \vec{e}_z\\
      1 & 0 & \pi/2\\
      \cos\Delta\varphi & \sin\Delta\varphi & \pi/2
    \end{vmatrix}.
\]

Expanding yields:
\begin{align*}
  n_x &= 0\cdot\tfrac{\pi}{2} - \tfrac{\pi}{2}\cdot\sin\Delta\varphi
       = -\frac{\pi}{2}\sin\Delta\varphi, \\
  n_y &= \tfrac{\pi}{2}\cdot\cos\Delta\varphi - 1\cdot\tfrac{\pi}{2}
       = \frac{\pi}{2}(\cos\Delta\varphi - 1), \\
  n_z &= 1\cdot\sin\Delta\varphi - 0\cdot\cos\Delta\varphi
       = \sin\Delta\varphi.
\end{align*}

For $\Delta\varphi = 120^\circ$ ($\sin 120^\circ = \tfrac{\sqrt{3}}{2}$,
$\cos 120^\circ = -\tfrac{1}{2}$):
\[
  \vec{n}
  = \begin{pmatrix}
      -\dfrac{\pi\sqrt{3}}{4}\\[6pt]
      -\dfrac{3\pi}{4}\\[6pt]
      \dfrac{\sqrt{3}}{2}
    \end{pmatrix}
  \approx
  \begin{pmatrix}-1.36034952\\-2.35619449\\0.86602540\end{pmatrix},
  \qquad
  |\vec{n}| \approx 2.85520635.
\]

The tilt angle of the triangular face relative to the $xy$-plane is the angle
between $\vec{n}$ and the $z$-axis:
\[
  \gamma
  = \arccos\!\left(\frac{|n_z|}{|\vec{n}|}\right)
  = \arccos\!\left(\frac{\sqrt{3}/2}{|\vec{n}|}\right)
  \approx \arccos(0.30331447).
\]

\begin{satz}[Tilt Angle of the Triangular Face]
\begin{equation}
  \gamma = \arccos\!\left(\frac{|{n}_z|}{|\vec{n}|}\right)
  \approx 72.34321285^\circ.
\end{equation}
The triangular face is therefore nearly perpendicular to the $xy$-plane,
consistent with the image of a steeply rising cone.
\end{satz}

% ─────────────────────────────────────────────────────────────────────────────
\section{Area of the Triangle at \texorpdfstring{$n = 19$}{n=19}}
% ─────────────────────────────────────────────────────────────────────────────

As a concrete example we compute the triangular area spanned by the
two generators of the twin zone at $n = 19$.

\subsection*{Step 1 --- Determine the Vertices}

The radius at $z = 19$:
\[
  r_{19}
  = \frac{2\cdot 19}{\pi} + \frac{3}{2\pi}
  = \frac{79}{2\pi}
  \approx 12.57324050.
\]

The three vertices of the triangle (apex $A$, and two points $B$, $C$
on the cone surface with $\Delta\varphi = 120^\circ$):
\begin{align*}
  A &= \left(0,\; 0,\; -\tfrac{3}{4}\right), \\
  B &= \left(r_{19},\; 0,\; 19\right)
     \approx (12.573241,\; 0,\; 19), \\
  C &= \left(r_{19}\cos 120^\circ,\; r_{19}\sin 120^\circ,\; 19\right)
     \approx (-6.286621,\; 10.888746,\; 19).
\end{align*}

\subsection*{Step 2 --- Determine the Side Lengths}

\textbf{Sides $|AB|$ and $|AC|$} (generators):

Since both generators lead to the same height $z = 19$,
$|AB| = |AC|$. The generator length is computed as:
\[
  |AB|
  = r_{19}\cdot|\vec{v}|
  = \frac{79}{2\pi}\cdot\sqrt{1+\frac{\pi^2}{4}}
  \approx 12.573241\cdot 1.862096
  \approx 23.41257946.
\]

\textbf{Base $|BC|$} (chord at radius $r_{19}$ and angle $120^\circ$):
\[
  |BC|
  = 2\,r_{19}\sin\!\left(\frac{120^\circ}{2}\right)
  = 2\,r_{19}\sin 60^\circ
  = 2\,r_{19}\cdot\frac{\sqrt{3}}{2}
  = r_{19}\sqrt{3}
  = \frac{79\sqrt{3}}{2\pi}
  \approx 21.77749137.
\]

\subsection*{Step 3 --- Area via Cross Product}

The edge vectors from point $A$:
\[
  \overrightarrow{AB} = B - A
  = \begin{pmatrix}r_{19}\\0\\19+\tfrac{3}{4}\end{pmatrix}
  = \begin{pmatrix}79/(2\pi)\\0\\79/4\end{pmatrix},
  \qquad
  \overrightarrow{AC} = C - A
  = \begin{pmatrix}-r_{19}/2\\ r_{19}\sqrt{3}/2\\79/4\end{pmatrix}.
\]

The cross product:
\[
  \overrightarrow{AB}\times\overrightarrow{AC}
  \approx (-215.052727,\; -372.482250,\; 136.906818),
  \qquad
  \left|\overrightarrow{AB}\times\overrightarrow{AC}\right|
  \approx 451.36922682.
\]

The area of the triangle:
\[
  F
  = \frac{1}{2}\left|\overrightarrow{AB}\times\overrightarrow{AC}\right|
  \approx \frac{451.369}{2}.
\]

\begin{satz}[Triangular Area at $n=19$]
\begin{equation}
  \boxed{F_{19} = \frac{1}{2}\left|\overrightarrow{AB}\times\overrightarrow{AC}\right|
  \approx 225.68461341}
\end{equation}
\end{satz}

\subsection*{Verification via Heron's Formula}

With side lengths
$a = |BC| \approx 21.77749137$,
$b = |AC| \approx 23.41257946$,
$c = |AB| \approx 23.41257946$ and
$s = \tfrac{a+b+c}{2} \approx 34.30132515$:
\[
  F = \sqrt{s(s-a)(s-b)(s-c)}
  \approx \sqrt{34.301\cdot 12.524\cdot 10.889\cdot 10.889}
  \approx 225.68461341 \quad\checkmark
\]

% ─────────────────────────────────────────────────────────────────────────────
\section{Summary of All Angles and Dimensions}
% ─────────────────────────────────────────────────────────────────────────────

\begin{table}[H]
\centering
\renewcommand{\arraystretch}{1.4}
\caption{Geometric characteristics of the cone generators.\protect\footnotemark}
\begin{tabular}{lll}
\toprule
Quantity & Exact Form & Numerical \\
\midrule
Cone slope & $dz/dr = \pi/2$ & $1.57079633$ \\
Generator length per $\Delta r=1$ & $\sqrt{1+\pi^2/4}$ & $1.86209589$ \\
Growth rate $dr/dn$ & $2/\pi$ & $0.63661977$ \\
$\Delta r$ minor step & $4/\pi$ & $1.27323954$ \\
$\Delta r$ major step & $8/\pi$ & $2.54647909$ \\
Azimuthal angle minor & $2\pi/3$ & $120^\circ$ \\
Azimuthal angle major & $4\pi/3$ & $240^\circ$ \\
Angle between generators (both) &
  $\arccos\!\bigl(\tfrac{\pi^2-2}{\pi^2+4}\bigr)$ & $55.43092771^\circ$ \\
Tilt angle of triangular face &
  $\arccos\!\bigl(\tfrac{|n_z|}{|\vec{n}|}\bigr)$ & $72.34321285^\circ$ \\
\midrule
Radius at $n=19$ & $79/(2\pi)$ & $12.57324050$ \\
Generator length at $n=19$ & $\tfrac{79}{2\pi}\sqrt{1+\pi^2/4}$ & $23.41257946$ \\
Base side $|BC|$ at $n=19$ & $\tfrac{79\sqrt{3}}{2\pi}$ & $21.77749137$ \\
Triangular area at $n=19$ & $\tfrac{1}{2}|\vec{AB}\times\vec{AC}|$ & $225.68461341$ \\
Half-opening angle $\alpha$ (to the axis) & $\arctan(2/\pi)$ & $32.48163659^\circ$ \\
Inclination angle $\beta$ (to the base) & $\arctan(\pi/2)$ & $57.51836341^\circ$ \\
Development angle $\varphi$ (unrolled mantle) & $4\pi/\sqrt{\pi^2+4}$ & $193.33053797^\circ$ \\
\bottomrule
\end{tabular}
\end{table}
\footnotetext{\label{rem:mantle_angles}%
Two angles complete the description of the cone. The inclination angle
$\beta = 90^\circ - \alpha = \arctan(\pi/2)$ is the angular form of the
cone slope $dz/dr = \pi/2$. Unrolling the lateral surface along a
generatrix produces a plane sector whose central angle --- the
development angle --- is $\varphi = 2\pi \sin\alpha =
4\pi/\sqrt{\pi^2+4} = 3.37425443~\mathrm{rad} = 193.33053797^\circ$, independent of
the height, as required for a cone. Two exact consequences of
$\tan\beta = \pi/2$: at every lattice point the height above the apex is
$h = z + \tfrac34 = Q/4$, while the circumference is $2\pi r = Q$
(Remark~\ref{rem:radius_arithmetic}). The apex height is therefore
exactly one quarter of the circumference --- for $n = 19$: circumference
$79$, apex height $79/4$ --- and the unrolled net of the cone is a
sector of $193.33^\circ$ whose arc through the lattice point $n$ has the
integer length $Q$.}
